\begin{table*}[htbp]
\centering\footnotesize
\begin{adjustbox}{max width=\textwidth}
\begin{tabular}{@{}lrrrrrrr@{}}
\toprule
Arm & Cases & Obj. recall (arm / Full) & Ans. cov. (arm / Full) & Coverage (arm / Full) & Unsupp. facts (arm / Full) & Cit. prec. (arm / Full) & Factor rec. (arm / Full) \\
\midrule
B0 & 50 & 0.00 / 0.70 (20) & 0.00 / 0.50 (20) & 0.86 / 0.83 (50) & 0.83 / 0.84 (43) & -- & 0.50 / 0.25 (8) \\
B1 & 50 & 0.00 / 0.70 (20) & 0.00 / 0.50 (20) & 0.90 / 0.83 (50) & 0.75 / 0.84 (45) & 0.00 / 0.71 (1) & 0.62 / 0.25 (8) \\
B4 & 50 & 0.93 / 0.70 (20) & 0.75 / 0.50 (20) & 0.45 / 0.83 (50) & 0.81 / 0.86 (50) & -- & 0.62 / 0.25 (8) \\
B6 & 50 & 0.00 / 0.70 (20) & 0.00 / 0.50 (20) & 1.00 / 0.83 (50) & 0.98 / 0.86 (50) & -- & 0.88 / 0.25 (8) \\
\bottomrule
\end{tabular}
\end{adjustbox}
\caption{Family v1: each baseline beside Full on the cases both answered without error. A cell is the arm's mean, then Full's mean on the same cases, over the cases where the metric is defined for both (the count in brackets). Unsupported facts: lower is better; the others: higher is better. A dash means the metric is undefined for the arm (for example, a baseline that cites no authority has no citation precision).}
\label{tab:main_v1}
\end{table*}